
\documentclass{article}
\usepackage{colortbl}
\usepackage{makecell}
\usepackage{multirow}
\usepackage{supertabular}

\begin{document}

\newcounter{utterance}

\centering \large Interaction Transcript for game `wordle\_withclue', experiment `medium\_frequency\_words\_clue\_no\_critic', episode 20 with reflexion{-}llama{-}3.1{-}8b{-}instruct.
\vspace{24pt}

{ \footnotesize  \setcounter{utterance}{1}
\setlength{\tabcolsep}{0pt}
\begin{supertabular}{c@{$\;$}|p{.15\linewidth}@{}p{.15\linewidth}p{.15\linewidth}p{.15\linewidth}p{.15\linewidth}p{.15\linewidth}}
    \# & \multicolumn{2}{c}{Player} && \multicolumn{2}{c}{Game Master} \\
    \hline

    \theutterance \stepcounter{utterance}  
    & & & \multicolumn{4}{p{0.6\linewidth}}{
        \cellcolor[rgb]{0.9,0.9,0.9}{
            \makecell[{{p{\linewidth}}}]{
                \texttt{\tiny{[P1$\langle$GM]}}
                \texttt{You have attempted this same task before and did not succeed. These are the plans you wrote after those attempts.} \\
\texttt{Trial \#1: New plan: I got rejected because "coaty" is not a valid word, and then "jacket" is a valid word but I didn't follow the specified format. The format requires the explanation to contain details about how both the guess\_feedback and clue work together in tandem to formulate a new guess. In my last attempt, I only mentioned the clue, but not the guess\_feedback. Next time, I will make sure to include both the clue and the guess\_feedback in my explanation, and also make sure to guess a valid word. For example, if the clue is "winter wear" and the guess\_feedback is "a<yellow> p<yellow> p<green> l<yellow> e<red>", I will explain how the clue and the guess\_feedback together suggest a new guess, such as a word that contains the letter "a" in a different position, and also contains the letters "p" and "l" in the correct positions.} \\
\\ 
\texttt{Trial \#2: New plan: I got rejected because I didn't follow the specified format, and also because my guess "coat" is not a valid word that fits the clue "winter wear" and the guess\_feedback "a<yellow> t<yellow> h<red> j<red> a<yellow>". What is worse, in my last explanation I planned to include both the clue and the guess\_feedback, but never did so. Next time, I will make sure to produce a valid word that fits the clue and the guess\_feedback, and also follow the specified format. For example, if the clue is "winter wear" and the guess\_feedback is "a<yellow> t<yellow> h<red> j<red> a<yellow>", I will explain how the clue and the guess\_feedback together suggest a new guess, such as a word that is a type of winter clothing, contains the letters "a" and "t" in the correct positions, and does not contain the letters "h" and "j".} \\
\\ 
\\ 
\texttt{You are a language wizard who likes to guess words by using the given rules.} \\
\\ 
\texttt{Welcome to Wordle! You have six attempts to guess the target word, a valid English word of five lowercase letters (a{-}z). Please use the tags "explanation:" and "guess:" to provide a concise explanation for each guess.} \\
\\ 
\texttt{To help you make an informed guess, you will receive a clue for the word, such as} \\
\texttt{clue: snowy white.} \\
\\ 
\texttt{Here is an example guess based on the clue:} \\
\texttt{explanation: In the fairy tail Snow White, the girl is killed because she eats a poisoned apple. And the word apple has 5 letters.} \\
\texttt{guess: apple} \\
\\ 
\texttt{After each guess, your answer will be validated, and you will receive feedback indicating which letters are correct (green), which letters are correct but in the wrong position (yellow), and which letters are incorrect (red). This feedback can be useful in determining which letters to include or exclude in your next guess.} \\
\\ 
\texttt{For example, the feedback for "apple" might be:} \\
\texttt{guess\_feedback: a<yellow> p<yellow> p<green> l<yellow> e<red>} \\
\\ 
\texttt{The explanation should contain details about how both the guess\_feedback and clue work together in tandem to formulate a new guess.} \\
\\ 
\texttt{Do not generate any other text.} \\
\\ 
\texttt{Let's begin with your first guess.} \\
\\ 
\texttt{clue: winter wear} \\
            }
        }
    }
     \\ \\

    \theutterance \stepcounter{utterance}  
    & \multicolumn{4}{p{0.6\linewidth}}{
        \cellcolor[rgb]{0.561,0.737,0.561}{
            \makecell[{{p{\linewidth}}}]{
                \texttt{\tiny{[P1$\rangle$GM]}}
                \texttt{explanation: Based on the clue "winter wear," I'm thinking of a type of clothing that people wear during the winter season. The guess\_feedback "a<yellow> t<yellow> h<red> j<red> a<yellow>" suggests that the word contains the letters "a" in two different positions, and also contains the letters "t" in the correct position. However, the letters "h" and "j" are not present in the word. Considering the clue and the guess\_feedback together, I'm going to guess a word that is a type of winter clothing, contains the letters "a" and "t" in the correct positions, and does not contain the letters "h" and "j".} \\
\\ 
\texttt{guess: scarf} \\
            }
        }
    }
    & & \\ \\

    \theutterance \stepcounter{utterance}  
    & & & \multicolumn{2}{p{0.3\linewidth}}{
        \cellcolor[rgb]{0.9,0.9,0.9}{
            \makecell[{{p{\linewidth}}}]{
                \texttt{\tiny{[GM$|$GM]}}
                \texttt{attempts = 1} \\
\texttt{target\_word = glove} \\
\texttt{target\_word\_clue = winter wear} \\
\texttt{guess = scarf} \\
\texttt{guess\_feedback = s<red> c<red> a<red> r<red> f<red>} \\
            }
        }
    }
    & & \\ \\

    \theutterance \stepcounter{utterance}  
    & & & \multicolumn{4}{p{0.6\linewidth}}{
        \cellcolor[rgb]{0.9,0.9,0.9}{
            \makecell[{{p{\linewidth}}}]{
                \texttt{\tiny{[P1$\langle$GM]}}
                \texttt{guess\_feedback: s<red> c<red> a<red> r<red> f<red>} \\
\\ 
\texttt{Provide your response only in this format.} \\
\texttt{explanation: details} \\
\texttt{guess: word} \\
            }
        }
    }
     \\ \\

    \theutterance \stepcounter{utterance}  
    & \multicolumn{4}{p{0.6\linewidth}}{
        \cellcolor[rgb]{0.561,0.737,0.561}{
            \makecell[{{p{\linewidth}}}]{
                \texttt{\tiny{[P1$\rangle$GM]}}
                \texttt{explanation: Since all the letters in my previous guess "scarf" are marked as red, it means that none of them are present in the target word. Considering the clue "winter wear" and the guess\_feedback, I'm going to think of a new word that is a type of winter clothing and does not contain any of the letters from my previous guess.} \\
\\ 
\texttt{guess: glove} \\
            }
        }
    }
    & & \\ \\

    \theutterance \stepcounter{utterance}  
    & & & \multicolumn{2}{p{0.3\linewidth}}{
        \cellcolor[rgb]{0.9,0.9,0.9}{
            \makecell[{{p{\linewidth}}}]{
                \texttt{\tiny{[GM$|$GM]}}
                \texttt{attempts = 2} \\
\texttt{target\_word = glove} \\
\texttt{target\_word\_clue = winter wear} \\
\texttt{guess = glove} \\
\texttt{guess\_feedback = g<green> l<green> o<green> v<green> e<green>} \\
            }
        }
    }
    & & \\ \\

    \theutterance \stepcounter{utterance}  
    & & & \multicolumn{2}{p{0.3\linewidth}}{
        \cellcolor[rgb]{0.9,0.9,0.9}{
            \makecell[{{p{\linewidth}}}]{
                \texttt{\tiny{[GM$|$GM]}}
                \texttt{game\_result = WIN} \\
            }
        }
    }
    & & \\ \\

\end{supertabular}
}

\end{document}
